\clearpage
\fancyhead[L]{\small GAUSS AN VON LINDENAU.}
\begin{center}
\textbf{GAUSS AN VON LINDENAU.}
\par\vspace{4pt}\hrule\vspace{4pt}
\fontsize{10}{12}\selectfont Monatliche Correspondenz. Band XX. S. 78, 79. 1809 Juli.
\par\vspace{4pt}\hrule\vspace{4pt}
\end{center}
\begin{flushright} Göttingen, den 14. August 1809.
\end{flushright}

\par\indent ---\enspace---\enspace--- Was ich Ihnen heute mitteilen kann, sind ein paar Beobachtungen der neuen Planeten, die ich in diesen Tagen gemacht habe. Mir kam es vor allem darauf an, Pallas rechtzeitig aufzusuchen. Letztes Jahr konnte ich die Elemente aus Mangel an guten \emph{vollständigen} Beobachtungen nicht verbessern (noch immer fehlen mir zuverlässigere \emph{Deklinationen}, und ich fürchte beinahe, dass ich außer meinen eigenen nur die mit dem Äquatorialsektor in Mailand gewonnenen werde verwenden können). Deshalb müssen wir uns diesmal mit der Ephemeride behelfen, die Herr Professor \textsc{Bode} im Jahrbuch für 1811 nach den 1807 bestimmten Elementen berechnet hat. --- Trotzdem hat es keine Mühe gekostet, den Planeten danach zu finden. Am 8. August hatte Herr Professor \textsc{Harding} die Gegend in den Fischen abgesucht, in der er stehen musste. Am 9. August erkannten wir ihn sofort in einem Stern 9. Größe, der sich merklich nach Süden bewegt hatte. Am 10. August war die Wiederauffindung völlig gewiss. Hier sind die Ergebnisse meiner Beobachtungen von diesen Tagen.


\begin{center}\fontsize{9}{11}\selectfont\setlength{\tabcolsep}{5pt}\renewcommand{\arraystretch}{1.15}\begin{tabular}{lrrr} 1809 & \shortstack{Mittlere Zeit\\in Göttingen.} & \shortstack{Scheinbare\\Rektaszension der $\PallasSymbol$} & \shortstack{Scheinbare nördliche\\Deklination.}\\\hline Aug. 9 & $11^{\text{u}}\;42^{\text{m}}\;48^{\text{s}}$ & $7^\circ\;51'\;44''$ & $2^\circ\;41'\;6''$\\ 10 & $11\;9\;58$ & $7\;50\;35$ & $2\;32\;37$\\
\end{tabular}\end{center}


\SourceLeaf{320}\par\indent Daraus geht hervor, dass die Ephemeride im astronomischen Jahrbuch die Deklinationen jetzt gut angibt, die Rektaszensionen aber um $20'$ zu klein.

\par\indent Am 10. August habe ich auch Ceres gefunden. Ein paar noch nicht vollständig reduzierte Beobachtungen zeigen zumindest, dass meine Ephemeride im Maiheft der M. C. die Position auf die Minute genau angibt. Ceres ist 8. bis 9. Größe.

\par\indent In der Nacht vom 12. auf den 13. glaubte Herr Professor \textsc{Harding}, auch Vesta bereits als Stern 9. bis 10. Größe erkannt zu haben, sehr nahe bei der Position, die die von Herrn Dr. \textsc{Schumacher} berechnete Ephemeride angibt. Das muss allerdings noch bestätigt werden. Dazu wird die erste klare Nacht genutzt werden.

